\section*{A.1 关键函数接口说明}
\addcontentsline{toc}{section}{A.1 关键函数接口说明}
\begin{table}[H]
\centering
\caption{关键函数说明}
\begin{tabular}{ll}
\toprule
函数名 & 作用 \\
\midrule
\texttt{simudata()} & 生成测角和测距仿真观测数据 \\
\texttt{TriRhoPosition()} & 三测站测距几何定位初值 \\
\texttt{objfun()} & 计算加权残差平方和 \\
\texttt{gradfun()} & 计算目标函数梯度 \\
\texttt{hessfun()} & 计算目标函数海森矩阵 \\
\texttt{plane\_position\_newton\_method\_detailed()} & 牛顿法主流程 \\
\texttt{plane\_position\_LM\_method\_detailed()} & 法方程迭代主流程 \\
\texttt{NewtonLaTeXOut()} & 输出牛顿法结果到 TeX 文档 \\
\texttt{LMLaTeXOut()} & 输出法方程结果到 TeX 文档 \\
\bottomrule
\end{tabular}
\end{table}

\section*{A.2 观测数据结构}
\addcontentsline{toc}{section}{A.2 观测数据结构}
\begin{itemize}
  \item 测角数组 \texttt{S1[100][6]}：前 3 列为测站坐标，第 4 列方位角，第 5 列俯仰角，第 6 列角度标准差；
  \item 测距数组 \texttt{S2[100][5]}：前 3 列为测站坐标，第 4 列斜距，第 5 列距离标准差。
\end{itemize}

\section*{A.3 迭代日志建议格式}
\addcontentsline{toc}{section}{A.3 迭代日志建议格式}
为便于调试和复现实验，建议每次迭代记录以下字段：
\[
\{k,\ F_k,\ \|\nabla F_k\|_2,\ \|\Delta\bm p_k\|_2,\ \bm p_k\}.
\]
对于法方程方法，建议额外记录
\[
\det(\bm A_k),\ \kappa(\bm A_k),
\]
用于判断病态程度。

\section*{A.4 导数校验方法}
\addcontentsline{toc}{section}{A.4 导数校验方法}
可采用中心差分检验解析梯度：
\[
\frac{\partial F}{\partial x_i}\approx
\frac{F(\bm p+h\bm e_i)-F(\bm p-h\bm e_i)}{2h},
\]
并计算相对误差
\[
\epsilon_i=\frac{|g_i^{ana}-g_i^{num}|}{\max(1,|g_i^{ana}|,|g_i^{num}|)}.
\]
若 $\epsilon_i<10^{-6}$，通常可认为导数实现正确。

\section*{A.5 伪代码补充}
\addcontentsline{toc}{section}{A.5 伪代码补充}
\begin{algorithm}[H]
\caption{导数一致性测试}
\begin{algorithmic}[1]
\State 给定测试点 $\bm p$ 和扰动 $h$
\For{$i=1$ to $3$}
\State 计算解析梯度分量 $g_i^{ana}$
\State 计算数值梯度分量 $g_i^{num}$
\State 计算相对误差 $\epsilon_i$
\If{$\epsilon_i$ 超阈值}
\State 标记该导数分量可能存在实现错误
\EndIf
\EndFor
\State 输出测试报告
\end{algorithmic}
\end{algorithm}

\section*{A.6 参数文件样例}
\addcontentsline{toc}{section}{A.6 参数文件样例}
建议采用文本或 JSON 保存实验参数，示例字段如下：
\begin{itemize}
  \item \texttt{angle\_sigma}：测角标准差（弧度）；
  \item \texttt{range\_sigma}：测距标准差（米）；
  \item \texttt{max\_iter}：最大迭代次数；
  \item \texttt{tol\_step}：步长收敛阈值；
  \item \texttt{seed}：随机种子；
  \item \texttt{station\_layout\_id}：测站布局编号。
\end{itemize}
